\chapter{Number-Line Representations and Quantitative Development}
\label{sec:p19-number-line-development}

Placing a written number on a blank line is not simply reading out a
quantity stored in the mind. The participant must understand the numeral,
the endpoints, the requested correspondence between numerical and spatial
intervals, and the response. An estimate therefore constrains an account
of those operations without uniquely displaying an internal code. This
distinction parallels the separation of geometric task, observed measure
and proposed transformation in \cref{sec:p19-mental-rotation}.
Here the empirical questions concern symbolic number-line estimation,
its relation to remembering numbers and learning arithmetic, and the
effects of a particular numerical board-game experience.

\section{Quantities, symbols and proposed magnitude codes}
\label{sec:p19-nl-quantities}

Seven visible dots and the written numeral 7 need not demand the same
operations. A dot comparison can succeed without knowing the conventional
symbol; knowing how to recite a counting sequence does not alone establish
which collection a spoken number denotes. Conversely, symbols can specify
an exact integer and support a written record even when an unaided visual
estimate of a large set is inexact. These are distinctions among tasks and
available information, not a claim that every nonsymbolic task is
approximate or that every response to a numeral is exact.
Small collections can be tracked as individual objects; larger collections
also permit approximate magnitude judgments.

The Opfer--Siegler synthesis treats noisy nonsymbolic magnitude processing
as a foundation for learning numerical symbols and proposes coexisting,
range-sensitive representations \cite[author draft, pp.~2--3,
10--15 and 21--28]{OpferSiegler2012QuantitativeDraftNL}.
One candidate is a compressed logarithmic magnitude code; another is a
linear magnitude with variability proportional to magnitude. Both can
produce ratio-sensitive discrimination. Thus discriminating 4 from 8 more
readily than 8 from 12 does not, on its own, choose between their mean-code
and noise assumptions. A linear code with constant absolute variability
is a third candidate, not the same model as a linear code with scalar
variability \cite[pp.~237--238]{SieglerOpfer2003EstimationNL}.

\begin{hbtheorem}{thm:p19-nl-ratio-identification}{proved=annex:p19-nl-ratio-identification}{A ratio effect does not identify the mean code}
Let \(\kappa>0\), and let the real noise \(\varepsilon\) have
mean zero, finite second moment, and
\(\mathbb E(e^{2\varepsilon/\kappa})<\infty\).
For independently drawn copies \(\varepsilon_x,\varepsilon_y\)
and positive quantities \(x,y\), define logarithmic comparison
codes \(L_x=\kappa\ln x+\varepsilon_x\).
Put \(m=\mathbb E(e^{\varepsilon/\kappa})\) and define
\(Z_x=e^{\varepsilon_x/\kappa}/m\), \(S_x=xZ_x\), with
the analogous definitions for \(y\).
Then \(m\) is finite and positive, the log codes have constant
variance, and the scalar codes have
\(\mathbb E(S_x)=x\) and
\(\operatorname{SD}(S_x)=x\operatorname{SD}(Z_x)\).
For every draw, \(L_x>L_y\) exactly when \(S_x>S_y\);
the comparison probability in both models depends only on \(x/y\).
Thus the same comparison observations can arise from a logarithmic
mean code with fixed additive noise or a linear mean code with
scalar multiplicative variability. This construction is a conditional
identification result, not an estimated distribution for infants
or evidence that either code explains every numerical task.
\end{hbtheorem}

The proposal must not erase the evidence modalities. Longer looking at a
changed infant display is a novelty response, not an explicit number-line
placement. In visual sets, numerosity can covary with area, contour,
density or object size. Varying these cues, or testing matching between
auditory sequences and visual collections, makes a purely visual-area
explanation less adequate; it does not establish every numerical concept
or eliminate duration and sequence alternatives automatically.
Object tracking is especially relevant to very small sets and occluded
objects. These controls are part of the explanatory problem, not details
that can be removed because a response is called numerical.

Similarly, an adult or child fMRI signal, infant scalp ERP, and a
trained macaque's single-cell response have different spatial resolution,
training histories and tasks. Anatomically overlapping activity is
consistent with related processing but does not establish identical codes,
one innate circuit or a causal explanation of a child's estimation error.
Comparisons among communities with different numerical vocabularies must
retain language, schooling, practice, translation and task familiarity:
performance in a sampled task does not rank the community's general
reasoning. The synthesis's cross-species, cultural and neural arguments
are broader than the school-age and preschool experiments below
\cite[author draft, pp.~6--15, 23--24 and 35--38]
{OpferSiegler2012QuantitativeDraftNL}.

\section{A bounded number-line task and its measurements}
\label{sec:p19-nl-task}

\begin{hbdefinition}{def:p19-nl-measures}{Position, accuracy and fit}
On a line with numerical endpoints 0 and \(U>0\), let \(x\) be the
presented number and \(y\) the numerical value obtained by converting
the participant's marked spatial position using the stipulated linear
scale. Both are in number units; the physical mark itself is a length.
Accurate placement has \(y=x\). Percentage absolute error is
\[
 \mathrm{PAE}=100\frac{|y-x|}{U}.
\]
An average of these errors retains trial inaccuracies. A fitted
relationship between estimates and targets answers a different question.
A best-fitting affine function \(y=a+bx\) can be highly linear without
having \(a=0\) and \(b=1\). Reported \(R^2\) values concern the
specified fit and data summaries, not the percentage of trials answered
exactly correctly. Group-median fits and averages of individual fit
statistics have different denominators and need not agree.
\end{hbdefinition}

For positive targets a fitted compressed mapping can use \(\ln x\),
where equal ratios give equal log differences. A zero endpoint on the
display does not make \(\ln0\) defined. A shifted-log function and an
unshifted positive-domain function are different models; any offset,
endpoint constraint or exclusion of zero must be specified. On the
reverse position-to-number task, an invertible logarithmic code predicts
an exponential relation only if the response implements that inverse.
Different tasks need not use identical decoding operations.

\section{Siegler and Opfer's two directions and two ranges}
\label{sec:p19-nl-siegler-opfer}

The 2003 study included 32 participants in each of four groups, total
\(N=128\): second graders, fourth graders, sixth graders and
undergraduates. Mean child ages were 7.9, 9.6 and 11.8 years
(standard deviations .4, .3 and .4). The children attended a suburban
school in an upper-middle-class area; the undergraduates were paid
volunteers. Those are the sampled populations, not a probability sample
of all children \cite[pp.~238--239]{SieglerOpfer2003EstimationNL}.

On the number-to-position task, NP, the participant marked where a
presented numeral belonged. On the position-to-number task, PN, the
participant supplied a numeral for a marked location. The 25-cm lines
were bounded by 0 and either 100 or 1,000. Two number sets were chosen
to distinguish logarithmic and linear functions and reduce simple
midpoint knowledge; for the 1,000 range one set included 4, 6, 18, 71,
230 and 780, the other 2, 6, 25, 86, 390 and 810. Each participant
was tested in two sessions one or two days apart. Task direction,
number set and four- versus thirty-second response limit were
counterbalanced and reversed in the second session; range order was
counterbalanced and items randomized. Range and available response
time were therefore experimentally varied within this design; age was
not randomly assigned.

For the 0--1,000 range the primary analysis fitted each age group's
median estimates. Second-graders' NP medians gave \(R^2=.95\) for
the logarithmic fit versus .63 for the linear fit
(\(t(10)=3.25\), \(p<.01\)); their PN medians gave .95 for the
exponential fit versus .62 for the linear fit
(\(t(10)=2.53\), \(p<.05\)). Fourth-graders' NP fits were .93
and .82, with no detected difference (\(t(10)=1.28\)); their PN
exponential and linear fits were .87 and .90. Sixth-graders' and
adults' NP medians had linear fits reported as 1.00, against
logarithmic fits .78 and .73. A rounded 1.00 fit is not a statement
of zero individual error. The PN linear fits were .97 and 1.00;
their advantages over exponential fits were reported only at
\(p<.10\), not at a conventional .05 threshold
\cite[pp.~239--240]{SieglerOpfer2003EstimationNL}.

Individual fits added a different comparison. The linear model best
fitted NP 1,000 estimates for 9\% of second graders, 38\% of fourth
graders, 72\% of sixth graders and 97\% of adults. On NP 100,
56\% of second graders and 69\% of fourth graders were best fitted
by a linear function. Thus some of the same children gave more linear
responses for the smaller range; a universal inability to follow the
instructions is insufficient as an explanation. This does not eliminate
range-sensitive endpoint comprehension or strategy choice.
For individual regressions involving zero responses, the article
explicitly added 1 to actual and estimated magnitudes to enable its
exponential regression calculation; this convention must not be
silently identified with every displayed logarithmic model
\cite[pp.~240--241]{SieglerOpfer2003EstimationNL}.

Variability also mattered. A landmark account proposed divisions at
250, 500 and 750: estimates nearer these references should be less
variable. It accounted for reported variance fractions .73 and .49
for sixth-graders' NP and PN variability, and .39 and .60 for adults'.
The candidate scalar-variability account explained less than 5\% of
NP variability in either of these groups. These are comparisons of
specified response models in a bounded task. They do not prove that
scalar variability never occurs in an internal magnitude representation,
and the later proportion account explicitly contests that inference.
The authors' proposed representations and landmark operations remain
explanations of the observations, not newly measured neural states
\cite[p.~242]{SieglerOpfer2003EstimationNL}.

\section{Proportion judgment as a competing construction}
\label{sec:p19-nl-proportion}

Barth and Paladino's alternative asks what a bounded task requires:
the target must be compared with a whole or with local reference
intervals, rather than estimated in isolation. Their model separates
bias in magnitude judgments from knowledge and use of reference
points \cite[pp.~125--127]{BarthPaladino2011ProportionNL}.
The later Slusser--Santiago--Barth study compares this account across
three age groups and pairs of numerical ranges
\cite[pp.~3--6 of the online-first article]
{SlusserEtAl2012OnlineProportionNL}.

\begin{hbdefinition}{def:p19-nl-proportion-model}{One cycle and a midpoint}
Let \(t\in[0,1]\) be a true proportion, and let \(\beta>0\)
be a power exponent. The proposed one-cycle proportion mapping is
\[
 P_\beta(t)=\frac{t^\beta}{t^\beta+(1-t)^\beta}.
\]
On a 0--\(U\) number line the corresponding predicted numerical
estimate is \(\widehat y=U P_\beta(x/U)\).
This specifies both endpoints without taking a logarithm of zero.
With a midpoint as an additional reference, the two-cycle version
applies the same mapping separately to each half:
\[
 Q_\beta(t)=
 \begin{cases}
   \tfrac12P_\beta(2t),&0\leq t\leq\tfrac12,\\
   \tfrac12+\tfrac12P_\beta(2t-1),&\tfrac12\leq t\leq1.
 \end{cases}
\]
The two formulas agree at the midpoint. These are predicted
response mappings under specified reference-point use, not claims
that a child's fitted exponent is directly a neural parameter.
An unbounded power mapping \(y=\alpha x^\beta\), with scaling
factor \(\alpha\), is a different candidate when an upper reference
is not effectively used. The free/fixed parameters and permitted
domains must be retained in comparisons.
\end{hbdefinition}

\begin{hbtheorem}{thm:p19-nl-proportion-bias}{proved=annex:p19-nl-proportion-bias}{Bias patterns of the stated proportional model}
The mapping \(P_\beta\) is continuous, strictly increasing,
fixes 0, \(1/2\) and 1, and satisfies
\(P_\beta(1-t)=1-P_\beta(t)\).
For \(0<\beta<1\), it overestimates proportions strictly
between 0 and \(1/2\), and underestimates those strictly between
\(1/2\) and 1. For \(\beta>1\) the inequalities reverse;
for \(\beta=1\), \(P_\beta(t)=t\).
The two-cycle \(Q_\beta\) fixes the quarter points as well as
the endpoints and midpoint. Its corresponding bias repeats in
each half: for \(0<\beta<1\), over, under, over, under on the
four successive open quarter intervals. Both maps are identical
to accurate linear placement when \(\beta=1\).
Thus accurate linear estimates alone cannot distinguish these
stipulated response constructions at that parameter value.
\end{hbtheorem}

This yields a discriminating observation beyond a generic statement
that younger children overestimate. A one-cycle model with
\(0<\beta<1\) predicts underestimation above the midpoint;
a two-cycle model predicts a renewed overestimate just above it.
Sampling mainly small targets can miss those differences. With
\(\beta=1/2\), for example, a one-cycle prediction at \(x=25\)
on a 0--100 line is \(100/(1+\sqrt3)\), about 36.6, whereas
the two-cycle prediction is exactly 25. These are model predictions,
not observations fabricated for a child.

Slusser and colleagues recruited 33 five- and six-year-olds,
24 seven- and eight-year-olds, and 30 eight- through ten-year-olds
from central Connecticut families, with some testing at local venues.
Socioeconomic status, race and ethnicity were not collected; community
census statistics cannot replace participant demographics.
The first experiment excluded eight children whose estimates failed
its rank-correlation condition and five with over 90\% of responses
in a tenth of the line, leaving 20. All 24 were retained in the second;
the third retained 28 after one response-pattern exclusion and one
parent-reported developmental-delay exclusion. These criteria select
interpretable response patterns and limit generalization to all
children initially recruited. The experiments are cross-sectional,
not observations of each child's developmental transition
\cite[pp.~5--10]{SlusserEtAl2012OnlineProportionNL}.

Each child received a 23-cm bounded line in a booklet and made a
mark for the printed, orally read numeral. A more familiar range
always preceded a less familiar one: 0--20 then 0--100, 0--100 then
0--1,000, or 0--1,000 then 0--100,000. A practice example explicitly
explained the midpoint of each range; the first experiment also
followed a PN task belonging to another study. Hence range familiarity,
order and prior practice are not fully separated. Target sampling was
spread across the range, unlike a deliberate emphasis on small targets.
For example, Experiment 2's 1,000 line used 24 targets from 3 through
997, including 475, 525 and 948. Appendix A gives the actual sets.

Models were compared with a small-sample-corrected Akaike criterion,
\(\mathrm{AIC_c}\), and leave-one-out prediction error, not merely
the largest \(R^2\). In the second experiment, Appendix C ranks
the two-cycle model first for both ranges. On 0--100 its reported
\(R^2=.990\), \(\mathrm{AIC_c}=54.994\), contrasts with the
linear model's .986 and 66.007. On 0--1,000 the corresponding
values are .989 and 167.550, against .970 and 192.574.
These are results for group medians under that candidate set, fitting
procedure and exclusions, not proof that no other account can work.
Figure 3's caption identifies the median for target 52 on the
1,000 line as an excluded statistical outlier: nearly half the
children placed it much too high, tending toward the location for
500. Such target-level exclusion is distinct from retaining all
24 participants. Its response error is itself part of the task's
numeral-interpretation problem, not evidence that those children
were making accurate proportion judgments on every target.
Holding out one target tests prediction of another target in these
records, not transfer to a newly sampled child or an arithmetic task
\cite[Appendices A and C]{SlusserEtAl2012OnlineProportionNL}.

Several printed numerical details require caution. Experiment 1's
results prose and Appendix B give .980 for the familiar-range
group's linear fit, whereas Figure 2A prints .937; Experiment 3's
body/figure and Appendix D also differ in some fit values. The
youngest group's 0--100 logarithmic fit is .838 in prose and .844
in Appendix B. These are not averaged or silently reconciled here.
The selected Appendix C comparisons above identify their own locator.
The paper also reports some older children's reversed, under-then-over
bias with \(\beta>1\), contrary to a simple progression restricted
to exponents below 1. Its suggested overcorrection explanation is
speculative, not a measured mechanism.

The proportion account preserves developmental change while
altering its interpretation: more usable endpoints or a midpoint,
less biased magnitude judgments, and better numeral knowledge can
change the response curve without a categorical replacement of
one underlying number code by another. A fitted preferred model
constrains the alternatives considered. It does not uniquely identify
the child's strategy, and the 2003 and later designs differ in ranges,
practice, sampling of targets and treatment of nonresponsive records.

\section{Remembering numbers is a separate outcome}
\label{sec:p19-nl-memory}

Thompson and Siegler tested memory for numerical content in short
stories, not memory for a spatial location or recall of a digit string.
After hearing a vignette containing numbers, children named four cartoon
characters and then answered questions about its quantities. In the
preschool tasks, recall was assessed by exact answers, by PAE using
20 as the range denominator, and by \(R^2_{\rm lin}\), the fit of
recalled numbers to presented numbers. A high linearity score can
coexist with systematically wrong recalls; exact recall and magnitude
error therefore remain separate outcomes
\cite[pp.~1276--1277]{ThompsonSiegler2010MemoryNL}.

Experiment 1 included 18 Pittsburgh-area prekindergartners, mean age
5.17 years (SD .40); 56\% were male and 94\% described as Caucasian.
The districts considered them at risk for school failure, mainly from
low income. Children marked targets 1--19 on a 20-cm, 0--20 line,
and categorized numbers into small, medium and big boxes. Six stories
used two small, medium or big numbers. Correlations between recall
linearity and the three magnitude measures were of magnitude .56
or .57. Exact recall was 39\% for small numbers, 6\% for medium
and 17\% for big numbers. The pattern is not a monotone decline
across all three observed categories, even though the reported
number-size correlation was negative.

Experiment 2 included 19 children from middle-income preschools,
mean age 4.96 (SD .52); 68\% were male, 68\% Caucasian and
32\% Asian American. A distributed set of twelve numbers from 1
through 19 replaced the clustered categories, and counting 20 chips
and identifying written numerals were added. One child answered 5 on
every recall trial, leaving only 18 for recall \(R^2_{\rm lin}\)
correlations. Number-line linearity correlated .53 with recall
linearity and -.47 with recall PAE; corresponding categorization
correlations were .66 and -.68. Counting and numeral-identification
correlations with recall did not meet the report's significance
threshold. This is not evidence of exactly zero relations, nor proof
that the selected correlations differ significantly from each other.
The changed population and target distribution provide a bounded
replication, not a removal of every shared language, attention or
instructional-history explanation
\cite[Tables 1--2]{ThompsonSiegler2010MemoryNL}.

Experiment 3 extended the comparison to 127 second graders, mean
age 8.22 (SD .44), 59\% male and 91\% Caucasian, from schools
near Pittsburgh; approximately 30\% qualified for subsidized lunches.
The three magnitude tasks used the same 22 positive targets between
2 and 938. Besides number-line placement on 0--1,000, children
drew a requested length in an unfamiliar unit, zips, with short and
long references, and filled an on-screen box with an estimated dot
quantity by holding increase/decrease controls, with 0 and 1,000-dot
references. These responses differ: a line mark, a drawn length and
a produced collection. Their shared bounded numerical structure
does not eliminate response or reference-point demands.
The recall stories now contained three numbers from a small,
medium or big set; the supplemental full story list is separate
from the main article's task description.

The authors predicted a representation-by-size interaction:
children whose magnitude estimates were best fitted by a linear
function should recall larger numbers more accurately, but not
necessarily smaller ones. They assumed both groups could use a
compressed representation for small values; that choice was not
independently observed during recall. With grouping by number-line
fit, recall PAE above 150 was 21\% for the linear-best-fit group
versus 32\% for the logarithmic-best-fit group
(\(F(1,125)=24.93\), \(p<.001\), reported \(d=.86\)). Below
150 the corresponding values were 7\% and 11\%, with
\(F(1,125)=3.72\), \(p>.05\). Similar contrasts occurred when
grouping by measurement or numerosity fit. The smaller-number
nonsignificance is not demonstrated equality. Exact recall of
hundreds digits, tens digits and ones digits was reported as 28\%,
15\% and 13\%, respectively, a pattern consistent with better
retention of coarse numerical gist
\cite[pp.~1278--1280 and Table 3]{ThompsonSiegler2010MemoryNL}.

These are associations between task measures, not a random assignment
to internal representations. The results constrain an explanation
that numerical memory depends only on verbatim numeral knowledge,
but the classification used to test the mechanism itself depends on
which estimation models were fitted. Shared knowledge of numerical
range, language comprehension or strategy selection can affect both
tasks. The study does not show that practising line marks causes
better memory for telephone numbers, nor that a low estimation
score diagnoses a general memory deficit.

\section{An arithmetic-learning intervention with four conditions}
\label{sec:p19-nl-arithmetic}

Booth and Siegler's study recruited 105 first graders in the second
half of the school year from four public schools in a predominantly
middle-income district. Mean age was 7.20 years (SD .40); there
were 57 boys and 48 girls and 96\% were reported Caucasian.
The report gives random assignment to child-generated displays
(26), computer-generated displays (25), both (27), or a control
with no magnitude display (27). Two children were absent at follow-up
and one did not complete the posttest number-line task because of
a computer malfunction; outcome analyses therefore need not have
the same denominator \cite[pp.~1019--1020]
{BoothSiegler2008ArithmeticNL}.

The pretest included 26 targets from 2 through 98 on a computerized
0--100 line, with extra sampling below 30; a mathematics achievement
test; forward digit span; and 13 addition problems. Children were
asked to retrieve an answer within six seconds rather than count
or calculate. Four preselected problems, \(9+18\), \(26+27\),
\(17+29\) and \(49+43\), were trained. These were fixed training
items, not each child's four easiest failures. Pretest accuracy
was only 7\% on these four. Restricting time and discouraging
calculation changes the question from unrestricted arithmetic
competence to learning and retrieving these answers.

Every condition encountered each trained problem four times in
each of three training sessions. The first two presentations involved
repeating the addends and answer; the next two involved reading them.
All were shown the correct sum during training, and each session
included a no-feedback test. Children met an experimenter individually
for three 10--15-minute sessions within one week and a five-minute
follow-up two weeks later. The thirteen-item test was repeated at
follow-up; the number-line task was repeated at the end of training
\cite[pp.~1020--1023 and Table 1]{BoothSiegler2008ArithmeticNL}.

In the computer condition, accurate red, blue and purple bars
represented the first addend, second addend and sum against a
0--100 line. A correct display made the two addend lengths join
end to end to give the sum length. For \(17+29\), their lengths
are 17\% and 29\% of the whole, and the joined length is 46\%.
Children in the child-generated condition estimated these magnitudes
themselves; those displays could be wrong. In the combined condition
their own estimates and the accurate comparison appeared on separate
lines. Controls saw the same arithmetic problems and answers but
no bars. The manipulation therefore supplies a particular accurate
visual relation, not simply more answer exposure. The combined
condition also introduces comparing and distinguishing correct from
self-generated information.

\begin{table}[htbp]
\centering
\small
\caption{Booth--Siegler Table 4: percent correct on four trained
addition problems. Parentheses are reported between-child standard
deviations, not standard errors or confidence intervals.}
\label{tab:p19-nl-arithmetic-results}
\begin{tabular}{lrrr}
\hline
Condition & Pretest & End of training & Two-week follow-up\\
\hline
Control & 6 (13) & 23 (24) & 10 (16)\\
Child generated & 4 (12) & 28 (25) & 13 (20)\\
Computer generated & 9 (18) & 41 (24) & 13 (22)\\
Child and computer & 10 (17) & 25 (28) & 14 (9)\\
\hline
\end{tabular}
\end{table}

The three-factor computer-present/child-present/time analysis
reported an interaction for correct answers
(\(F(1,101)=5.01\), \(p<.05\), \(\eta^2=.05\)). The isolated
computer condition improved by 32 percentage points, compared
with 17, 24 and 15 in the control, child and combined conditions.
Reported least-significant-difference comparisons at training end
favoured computer-only over control and combined, with
\(p<.05\); no unreported multiplicity adjustment is supplied.
At the two-week comparison, only a time main effect was detected
for percent correct; no instructional-condition advantage was
detected. The PAE follow-up analysis likewise had no significant
main effects or interactions. Thus the immediate learning
contrast cannot be relabelled as a demonstrated durable advantage
\cite[pp.~1025--1026 and Table 4]{BoothSiegler2008ArithmeticNL}.

Number-line PAE also changed: children exposed to accurate computer
displays improved from 10\% to 9\%, whereas those not exposed
went from 11\% to 12\%, with an exposure-by-time interaction
\(F(1,100)=4.87\), \(p<.05\), \(\eta^2=.05\).
Baseline estimation linearity predicted arithmetic learning even
after the selected achievement, prior-addition and digit-span
covariates. This predictive relation is observational: statistical
adjustment does not randomly assign the predictor or establish that
all confounding has disappeared
\cite[pp.~1024--1027]{BoothSiegler2008ArithmeticNL}.

Random assignment supports a causal contrast between these
instructional packages in this setting, subject to the reported
missing outcomes. It does not by itself establish a linear internal
code as the mediator. Correct visual answer-size constraints,
attention, comparative encoding and source-monitoring demands are
also changed. The authors' speculation that the combined display
overloaded source monitoring was not independently tested.
The results specifically fail to support the prediction that requiring
self-generated representations would add to the accurate display's
benefit; they are not a theorem that self-generation or manipulatives
are always ineffective.

A compressed code alone does not prescribe an arithmetic algorithm.
If a process simply adds \(\ln17\) and \(\ln29\) and decodes
by exponentiation, it obtains \(17\cdot29=493\), not 46.
The article discusses adding encoded values as a possible reason
for poor approximations; a psychological operation and decoding
rule are additional premises, not consequences of saying that a
representation is logarithmic. Correct addition on a physical
linear display and an empirical hypothesis about mental operations
must remain distinct.

\section{Numerical board games: the intervention and its limits}
\label{sec:p19-nl-board-games}

Siegler and Ramani's 2008 \emph{Developmental Science} article
first compared 58 preschoolers aged 4.0--5.1 years: 36 from
low-income urban settings and 22 from a predominantly
upper-middle-class university preschool. Almost all families in the
former group received childcare subsidies. These observational
groups differ in more than income, and income was not an
experimental treatment. On a 25-cm, 0--10 line each target 1--10
was estimated twice, with no corrective feedback. The groups'
median estimates had linear \(R^2\) values .66 and .94;
mean individual fits were .15 and .60. These are two different
summary statistics, not conflicting estimates of one quantity
\cite[pp.~657--658]{SieglerRamani2008BoardNL}.

The same 36 low-income children entered Experiment 2, randomly
assigned 18 per condition, with mean ages 4.6 and 4.7 years.
Both Great Race boards were 50 by 30 cm, with eleven equal
horizontal squares including Start, and rabbit or bear tokens.
The numerical board labelled the ten further squares 1--10;
the other used colours. Children played individually with an
experimenter for four fifteen-minute sessions within two weeks,
approximately thirty games lasting two to four minutes each.
The numerical spinner indicated 1 or 2; children moved that
many squares while saying the numbers on the visited squares.
For example, moving from 3 by two meant saying 4, then 5,
not restarting a count at 1. The colour game's spinner and spoken
labels referred to colours. Errors received an adult naming
correction followed by the child's repetition during movement.
The number-line assessment occurred at the start of session one
and end of session four \cite[p.~658]{SieglerRamani2008BoardNL}.

The outcome was not an exactly correct answer count alone.
The numerical group's mean individual linear \(R^2\) increased
from .15 to .61; the colour group's mean stayed .18.
The posttest between-condition comparison was
\(t(34)=4.85\), \(p<.001\), reported \(d=1.62\).
Mean PAE fell from 28\% to 20\% in the numerical group
(\(t(17)=2.43\), \(p<.05\), \(d=.71\)), while the colour
group went from 27\% to 28\%. Mean fitted slopes went from
.23 to .71 in the numerical condition and .30 to .24 in the
colour condition. The percentage of pairwise estimates correctly
ordered increased from 62\% to 81\% in the numerical group;
the control remained approximately 62\%. Pairwise order is not
exact placement: a child can order all numbers correctly while
compressing every distance \cite[pp.~658--659]
{SieglerRamani2008BoardNL}.

The article's later catch-up paragraph calls numerical-game
posttest performance equivalent to the higher-income comparison
group, but this was not an equivalence-test design. Moreover it
prints individual slopes .65 and .66 there, whereas its detailed
results give .71 for the intervention group's posttest and .70
for the higher-income group's baseline. The matched fit values
and 81\% ordered-pair figures do not repair the slope discrepancy.
Neither a nonsignificant comparison nor this paragraph establishes
that all quantitative disadvantages disappeared.

The causal inference supported by random assignment is about
this numerical-game package compared with the colour package
over the observed period. Several cues covary in the numerical
game: location, accumulated moves, elapsed time, spoken and heard
number names, counting onward and feedback. These correlated
cues motivate learning a spatial-numerical mapping; the design
does not isolate which cue is necessary, whether a compressed
internal scale was replaced, or whether a different task strategy
was learned. The participant-blindness and assessor-blindness
conditions are not established by the report.

The article's discussion cites a separate \emph{Child Development}
study for counting, numeral identification, magnitude-comparison
and nine-week retention results. Those are not measured endpoints
of this experiment. An immediate number-line effect therefore
cannot be promoted here to durable broad mathematical transfer.
The present study likewise does not compare linear with circular
boards: that is a distinct later experiment. General classroom
recommendations require additional populations, implementation,
duration and outcome evidence, not merely an income-group
correlation or the low cost of the activity.

\section{The inferential work done by each comparison}
\label{sec:p19-nl-warrants}

\begin{hbmethod}{def:p19-nl-warrant-record}{Identify the response before the mechanism}
State the numerical range and endpoint knowledge, symbolic or
nonsymbolic input, spatial or verbal response, trial sampling and
practice. Keep participant counts distinct from targets, trials
and median data points. Report exact accuracy, signed or absolute
error, response variability and model fit with their actual
denominators. State which parameters were estimated, fixed or
constrained, which candidate models were compared and which records
were excluded. For an intervention, distinguish an immediate
trained-task contrast, an untrained task and delayed retention.
Then identify which process hypotheses predict different outcomes,
and which can reproduce the same observed mapping.
\end{hbmethod}

In these records, context-sensitive NP/PN estimates constrain a
single, context-invariant response scale; cyclic bias constrains a
logarithmic-versus-linear-only candidate set; numerical-story recall
adds an outcome beyond line placement; randomized visual displays
and games support bounded intervention contrasts. None alone shows
the complete route from nonsymbolic discrimination to symbolic
arithmetic, establishes a neural substrate, or turns an observed age
association into a universal developmental timetable.

Failure cases make the distinctions concrete. Perfectly ordered
marks can have a wrong slope; a curve close to a straight line can
retain systematic local bias; excluding uninterpretable younger
children can improve model fits while reducing the population to
which those fits apply. Better immediate arithmetic retrieval can
lose its between-condition advantage at follow-up. Calling each
outcome number sense conceals these different inferential jobs.
The proportional-model result above is a mathematical conditional;
the experiments assess how well such constructions describe their
recorded responses. Neither supplies a normative rule that an older
child's faster or more linear-looking response is always the better
reasoning process.
